\section{Gestión de la Información (Cumplimiento ISO 26515)}
\label{sec:s1_gestion}

El incremento de seguridad se acompaña de su documentación técnica, generada en paralelo al
código conforme a ISO/IEC/IEEE 26515. Se publica el primer borrador de la especificación de la
API (Sección~\ref{ssec:s1_gi_openapi}) y el diagrama de componentes del módulo de seguridad
(Sección~\ref{ssec:s1_gi_diagrama}).\par

\subsection{Primer Borrador Formal de la Especificación de la API}
\label{ssec:s1_gi_openapi}
Se publicó el primer borrador de la especificación OpenAPI mediante \comando{springdoc-openapi},
accesible en \comando{/swagger-ui.html} y \comando{/v3/api-docs}. La especificación documenta los
endpoints de \comando{auth} y \comando{security}, sus DTO validados (Jakarta Bean Validation) y
los códigos de respuesta envueltos en \comando{ApiResponse}.

\begin{itemize}[noitemsep]
    \item \textbf{Contrato vivo:} la spec se regenera en cada \textit{build}, evitando desfase
    con el código.
    \item \textbf{Ejemplos de petición/respuesta:} documentados por endpoint para facilitar el
    consumo desde el frontend.
    \item \textbf{Códigos normalizados:} 400 (validación), 401 (no autenticado), 403 (rol
    insuficiente) y 200/201 en éxito.
\end{itemize}

\subsection{Diagrama de Componentes del Módulo de Seguridad}
\label{ssec:s1_gi_diagrama}
Se elaboró un diagrama de componentes centrado en el flujo de seguridad (Figura~\ref{fig:s1_componentes}):
el cliente atraviesa la cadena de filtros hasta el \textit{Resource Server}, que valida el JWT
contra el JWKS de Supabase antes de alcanzar los controladores de dominio.

\begin{figure}[!ht]
    \centering
    \begin{tikzpicture}[
        node distance=0.7cm and 1.0cm,
        comp/.style={draw=colorPrimario, fill=headerTabla, thick, rounded corners=2pt,
            text width=2.6cm, align=center, minimum height=1.0cm, font=\sffamily\scriptsize},
        ext/.style={draw=grisISO, fill=grisTecnico, thick, rounded corners=2pt,
            text width=2.4cm, align=center, minimum height=1.0cm, font=\sffamily\scriptsize},
        rel/.style={-{Latex[length=2mm]}, colorPrimario, font=\sffamily\tiny}
    ]
        \node[comp] (cli) {Cliente (SPA)};
        \node[comp, right=of cli] (chain) {Cadena de filtros};
        \node[comp, right=of chain] (rs) {Resource Server};
        \node[ext, below=of rs] (jwks) {JWKS Supabase};
        \node[comp, left=of jwks] (ctrl) {Controladores\\auth / security};

        \draw[rel] (cli) -- (chain);
        \draw[rel] (chain) -- (rs);
        \draw[rel] (rs) -- node[right]{valida} (jwks);
        \draw[rel] (rs) -- node[above]{ok} (ctrl);
    \end{tikzpicture}
    \caption{Diagrama de componentes del módulo de seguridad (Sprint 1).}
    \label{fig:s1_componentes}
\end{figure}

La tabla Historias~vs.~Documentación se actualizó marcando la EPIC-01 como documentada, cerrando
el ciclo de \textit{Continuous Documentation} del sprint.

\begin{TablaHSDoc}
    EPIC-01 & Autenticación y cuentas & Sí (S1) & Sí (S1) \\ \hline
\end{TablaHSDoc}

\subsection{Productos de Información del Incremento (ISO 15289)}
\label{ssec:s1_gi_productos}
La tabla~\ref{tab:s1_productos} relaciona cada producto de información generado con su artefacto
de origen en el código, conforme a ISO/IEC/IEEE 15289.

\begin{table}[!ht]
    \centering
    \renewcommand{\arraystretch}{1.3}
    \fontsize{9}{10.8}\selectfont\sffamily
    \begin{tabularx}{\textwidth}{|X|X|}
        \hline
        \rowcolor{colorPrimario}
        \textcolor{white}{\textbf{Producto de información}} & \textcolor{white}{\textbf{Artefacto de origen}} \\ \hline
        Especificación OpenAPI (auth/security) & \comando{springdoc-openapi} + controladores. \\ \hline
        \rowcolor{grisTecnico}
        Diagrama de componentes de seguridad & Cadena de filtros de \comando{config/}. \\ \hline
        Guía del flujo OTP & \comando{SecurityController} + \comando{profile\_security\_codes}. \\ \hline
    \end{tabularx}
    \caption{Productos de información del Sprint 1 y su trazabilidad.}
    \label{tab:s1_productos}
\end{table}

\clearpage
